% !PRINTED 387
\weberpage{387}

%<<CONT
so hat $b$ einen der Werthe 1, 2, 3 \ldots $\lambda$, jedoch mit der Beschränkung, dass bei geradem $D$ nur die geraden, bei ungeradem $D$ nur die ungeraden unter diesen Zahlen für $b$ zu nehmen sind. Es ist dann $a$ und $c$ so zu bestimmen, dass

\weq{6}{\frac{D-b^2}{4} = ac,}

und dass

\weq{7}{\frac{\lambda-b+1}{2} \leqq c \leqq \frac{\lambda+b}{2},}

und $a$ muss zwischen denselben Grenzen liegen. Nur einer von diesen beiden Grenzwerthen ist eine ganze Zahl, der andere kann für die untere Grenze durch die nächst grössere, für die obere Grenze durch die nächst kleinere ganze Zahl ersetzt werden. Es ist also für $b$ nur eine endliche Zahl von Werthen zulässig; dann ist $(D-b^2):4$ nur auf eine endliche Anzahl von Arten in zwei Factoren zerlegbar, und von diesen Zerlegungen sind nur die beizubehalten, in denen beide Factoren der Bedingung (7) genügen. Ausserdem sind noch solche Combinationen wegzulassen, in denen $a$, $b$, $c$ einen gemeinsamen Theiler bekommen. Darin liegt das Mittel, um für ein gegebenes $D$ alle reducirten Zahlen wirklich zu bestimmen.

Wir wollen für das Folgende die so bestimmten Zahlen $\omega$ durch das Symbol

\weq{8}{\omega = \frac{\sqrt{D}+b}{2c} = \frac{2a}{\sqrt{D}-b} = \{a, b, c\}}

bezeichnen.

Aus 3. \S.\,124 und 4. \S.\,125 ergiebt sich nun der folgende wichtige Satz:

\gesperrt{5. Die Kettenbruchentwickelung einer reducirten Zahl $\omega$ ist periodisch,}

d. h. wenn $\omega$ in einen Kettenbruch

\weq{9}{\omega = (\alpha_0, \alpha_1, \alpha_2 \ldots)}

entwickelt wird, so kehren immer nach einer bestimmten Anzahl von Theilnennern dieselben Theilnenner in derselben Reihenfolge wieder, also, wenn die Periode aus $\nu$ Gliedern besteht, so ist
\[
\begin{array}{ccccc}
\alpha_0 &=& \alpha_\nu &=& \alpha_{2\nu} \ldots\\
\alpha_1 &=& \alpha_{\nu+1} &=& \alpha_{2\nu+1} \ldots\\
\multicolumn{5}{c}{\cdots\cdots\cdots\cdots\cdots\cdots\cdots\cdots\cdots}\\
\alpha_{\nu-1} &=& \alpha_{2\nu-1} &=& \alpha_{3\nu-1} \ldots,
\end{array}
\]
%CONT>>
